\documentclass[a4paper,11pt]{article}

\usepackage[utf8]{inputenc}
\usepackage[T1]{fontenc}
\usepackage[brazilian]{babel}
\usepackage[margin=2.5cm]{geometry}
\usepackage{graphicx}
\usepackage{booktabs}
\usepackage{amsmath,amssymb}
\usepackage{listings}
\usepackage{xcolor}
\usepackage{hyperref}

\hypersetup{colorlinks=true, linkcolor=blue!60!black, urlcolor=blue!60!black, citecolor=blue!60!black}

\definecolor{codegray}{rgb}{0.95,0.95,0.95}
\lstset{
    basicstyle=\ttfamily\footnotesize,
    backgroundcolor=\color{codegray},
    breaklines=true,
    frame=single,
    framerule=0pt,
    xleftmargin=0.5em,
    xrightmargin=0.5em,
    showstringspaces=false,
    keywordstyle=\color{blue!60!black}\bfseries,
    commentstyle=\color{green!40!black}\itshape,
    stringstyle=\color{red!50!black},
    language=Java
}

\title{TP01 -- Pontes e Caminho Euleriano\\
\large Métodos Naïve e Tarjan + Algoritmo de Fleury}
\author{PUC Minas -- Ciência da Computação\\
Teoria dos Grafos e Computabilidade\\
Prof. Zenilton Kleber Gonçalves do Patrocínio Júnior}
\date{\today}

\begin{document}

\maketitle

\section{Introdução}

Uma ponte é uma aresta cuja remoção desconecta o grafo. O conceito aparece em vários problemas, e o que motivou esse trabalho foi o algoritmo de Fleury, que constrói caminhos eulerianos -- sequências que percorrem cada aresta do grafo exatamente uma vez. A cada passo o Fleury precisa saber se uma aresta candidata é ponte; portanto, a estratégia escolhida para detectar pontes muda o comportamento prático do algoritmo.

Aqui implementamos os dois métodos pedidos pelo enunciado, o naïve (que testa conectividade após remover cada aresta) e o de Tarjan (1974), que faz tudo em uma única DFS calculando \textit{lowlink}. Os dois são acoplados ao Fleury via uma interface, o que permite trocar a estratégia sem mexer no resto. Em seguida, comparamos tempo médio em grafos aleatórios de 100, 1.000, 10.000 e 100.000 vértices, nos três tipos eulerianos.

\section{Fundamentos}

\subsection{Quando existe caminho ou circuito euleriano}

Pelo teorema de Euler, num grafo conexo:
\begin{itemize}
    \item existe circuito euleriano sse todos os vértices têm grau par;
    \item existe caminho euleriano (sem ser circuito) sse exatamente dois vértices têm grau ímpar;
    \item caso contrário, nem caminho nem circuito existem.
\end{itemize}

Essa classificação é a primeira coisa que o Fleury faz. Se o grafo cai no terceiro caso, abortamos antes de tentar construir o caminho.

\subsection{Método naïve}

O naïve é literalmente a definição de ponte virada em código. Para cada aresta $(u,v)$, removemos, perguntamos se $v$ ainda é alcançável a partir de $u$ por DFS, e reinserimos. Se a resposta foi não, $(u,v)$ era ponte. O custo é $O(E \cdot (V+E))$, uma DFS por aresta.

Uma otimização que usamos: em vez de \texttt{isConnected()} (que varre o grafo todo), o naïve usa \texttt{isReachable(u, v)}, que para a DFS assim que encontra $v$. Em grafos esparsos, essa terminação antecipada faz uma diferença grande, e isso vai reaparecer na análise dos experimentos.

\subsection{Método de Tarjan}

A ideia do Tarjan é trocar $E$ DFS por uma só. Em uma única passada, mantemos para cada vértice $v$:
\begin{itemize}
    \item $\mathrm{disc}(v)$, o instante em que $v$ foi descoberto;
    \item $\mathrm{low}(v)$, o menor $\mathrm{disc}$ alcançável a partir de $v$ usando arestas da árvore DFS e no máximo uma aresta de retorno.
\end{itemize}

Quando $v$ é filho de $u$ na árvore DFS, a aresta $(u,v)$ é ponte se e somente se $\mathrm{low}(v) > \mathrm{disc}(u)$. Em palavras: nada que sai da subárvore de $v$ chega em $u$ ou em algum ancestral de $u$ por outro caminho. Custo $O(V+E)$.

Implementamos a versão iterativa, com pilha explícita. A versão recursiva existe no código (\texttt{TarjanBridgeFinderRec.java}) só por valor didático; em $V \geq 10.000$ ela estoura o stack do Java.

\subsection{Fleury}

O algoritmo, na prática, é:
\begin{enumerate}
    \item classifique o grafo; se for não-euleriano, aborte;
    \item escolha $v_0$: ímpar (semi-euleriano) ou qualquer vértice de grau positivo (euleriano);
    \item estando em $u$, prefira sempre uma aresta que não seja ponte. Pontes são último recurso;
    \item remova a aresta percorrida, mova-se para o vizinho, repita até esgotar arestas.
\end{enumerate}

A regra de evitar pontes é o que evita ficar preso. Atravessar uma ponte sem necessidade pode isolar arestas que ainda precisariam ser percorridas, e aí o algoritmo trava antes do fim.

\section{Implementação}

Tudo em Java 21. Os arquivos estão organizados assim:

\begin{table}[h]
\centering
\caption{Arquivos do projeto}
\small
\begin{tabular}{ll}
\toprule
Arquivo & O que faz \\
\midrule
\texttt{Graph.java}                 & Lista de adjacência, contador de versão e de arestas \\
\texttt{GraphUtils.java}            & DFS, BFS, conectividade, classificação euleriana \\
\texttt{GraphGenerator.java}        & Gera grafos aleatórios dos três tipos \\
\texttt{BridgeFinder.java}          & Interface (\textit{Strategy}) para detectar pontes \\
\texttt{NaiveBridgeFinder.java}     & Naïve: remove + testa + reinsere \\
\texttt{TarjanBridgeFinder.java}    & Tarjan iterativo com cache versionado \\
\texttt{TarjanBridgeFinderRec.java} & Tarjan recursivo (didático, não usado) \\
\texttt{FleuryEulerPath.java}       & Fleury parametrizado por \texttt{BridgeFinder} \\
\texttt{Experiments.java}           & Bateria de experimentos $\rightarrow$ CSV + tabelas \\
\texttt{Main.java}                  & Menu interativo + CLI batch \\
\bottomrule
\end{tabular}
\end{table}

\subsection{Estrutura de dados}

Lista de adjacência. A alternativa, matriz, ocuparia $O(V^2)$, o que para $V=100.000$ daria $10^{10}$ posições -- impossível na memória disponível. Lista de adjacência custa $O(V+E)$ e ainda permite remoção/inserção de arestas em tempo razoável, que é o que o naïve precisa fazer o tempo todo.

\subsection{Interface comum para pontes}

A interface \texttt{BridgeFinder} é o ponto de troca de estratégia. O Fleury não sabe se está usando naïve ou Tarjan, só pergunta se a aresta é ponte:

\begin{lstlisting}
public interface BridgeFinder {
    List<int[]> findBridges(Graph g);
    boolean     isBridge(Graph g, int u, int v);
    String      getName();
}
\end{lstlisting}

\subsection{Cache versionado no Tarjan}

O \texttt{isBridge} do Tarjan internamente chama \texttt{findBridges} (custa $O(V+E)$) e procura a aresta no resultado. O Fleury chama \texttt{isBridge} várias vezes seguidas no mesmo estado do grafo, então faz sentido guardar o resultado.

A solução: um contador de versão no \texttt{Graph}, incrementado a cada \texttt{addEdge}/\texttt{removeEdge}. O Tarjan armazena (versão observada, conjunto de pontes) e só recomputa se a versão atual diferir. Sem isso, Fleury+Tarjan seria $O(\mathrm{deg}\cdot E\cdot(V+E))$; com cache, $O(E\cdot(V+E))$.

\section{Experimentos}

\subsection{Configuração}

Linux x86\_64, JDK 21, JVM com \texttt{-Xss64m -Xmx2g}. Cada configuração roda 3 vezes; reportamos a média. Tamanhos $\{100,\ 1.000,\ 10.000,\ 100.000\}$ e densidade alvo de aproximadamente $1{,}1\,V$ arestas (grafo esparso, próximo de uma árvore com algumas arestas extras).

\subsection{Tempo médio para encontrar todas as pontes}

\begin{table}[h]
\centering
\caption{Tempo médio para encontrar todas as pontes do grafo (\texttt{findBridges}). ``--'' significa que a configuração foi pulada por inviabilidade de tempo.}
\label{tab:bridges}
\begin{tabular}{rlrrr}
\toprule
$N$ & Tipo & Arestas & Naïve (ms) & Tarjan (ms) \\
\midrule
$100$       & Eulerian      & $118$    & $4{,}59$    & $0{,}24$ \\
$100$       & Semi          & $121$    & $1{,}50$    & $0{,}08$ \\
$100$       & Não-Euler.    & $118$    & $1{,}45$    & $0{,}08$ \\
\midrule
$1.000$     & Eulerian      & $1.088$  & $82{,}40$   & $0{,}76$ \\
$1.000$     & Semi          & $1.093$  & $43{,}32$   & $0{,}65$ \\
$1.000$     & Não-Euler.    & $1.092$  & $44{,}55$   & $0{,}63$ \\
\midrule
$10.000$    & Eulerian      & $10.909$ & --          & $1{,}22$ \\
$10.000$    & Semi          & $10.884$ & --          & $1{,}16$ \\
$10.000$    & Não-Euler.    & $10.916$ & --          & $1{,}24$ \\
\midrule
$100.000$   & Eulerian      & $108.949$ & --         & $19{,}87$ \\
$100.000$   & Semi          & $109.011$ & --         & $19{,}07$ \\
$100.000$   & Não-Euler.    & $108.982$ & --         & $21{,}17$ \\
\bottomrule
\end{tabular}
\end{table}

O resultado é o esperado pela teoria: Tarjan ganha em todos os casos. Em $N=100$ a diferença é de cerca de 19$\times$; em $N=1.000$ chega a 108$\times$ no caso eulerian (82{,}4~ms contra 0{,}76~ms). Em $N=10.000$ e $100.000$ o naïve foi pulado porque a estimativa passa de cinco minutos por execução. Vale notar que o Tarjan resolve o caso de cem mil vértices em cerca de 20~ms.

\subsection{Tempo médio do Fleury completo}

\begin{table}[h]
\centering
\caption{Tempo médio para encontrar o caminho euleriano via Fleury. Em grafos não-eulerianos, o algoritmo aborta na classificação inicial.}
\label{tab:fleury}
\begin{tabular}{rlrrr}
\toprule
$N$ & Tipo & Arestas & Fleury+Naïve (ms) & Fleury+Tarjan (ms) \\
\midrule
$100$       & Eulerian      & $118$    & $2{,}28$  & $0{,}91$ \\
$100$       & Semi          & $121$    & $0{,}53$  & $0{,}66$ \\
$100$       & Não-Euler.    & $118$    & $0{,}08$  & $0{,}07$ \\
\midrule
$1.000$     & Eulerian      & $1.088$  & $5{,}23$  & $28{,}39$ \\
$1.000$     & Semi          & $1.093$  & $5{,}61$  & $29{,}62$ \\
$1.000$     & Não-Euler.    & $1.092$  & $0{,}58$  & $0{,}55$ \\
\midrule
$10.000$    & Eulerian      & $10.909$ & --        & $20.920{,}13$ \\
$10.000$    & Semi          & $10.884$ & --        & $20.526{,}37$ \\
$10.000$    & Não-Euler.    & $10.916$ & --        & $1{,}65$ \\
\midrule
$100.000$   & Eulerian      & $108.949$ & --       & --$^*$ \\
$100.000$   & Semi          & $109.011$ & --       & --$^*$ \\
$100.000$   & Não-Euler.    & $108.982$ & --       & --$^*$ \\
\bottomrule
\end{tabular}\\[0.4em]
{\small $^*$ Em $N=100.000$ a estimativa passa de 30~min por execução: $E\cdot(V+E)\approx 2{,}2\times 10^{10}$ operações elementares, o que dá cerca de 18~min com $\sim 50$~ns por operação, e a tabela exigiria três repetições por tipo.}
\end{table}

A leitura da tabela tem alguns pontos interessantes:

Para grafos não-eulerianos, os dois métodos terminam em menos de 1~ms. Faz sentido: a classificação detecta na hora que não existe caminho e o Fleury aborta sem percorrer nada.

Em $N=100$, Fleury+Tarjan é mais rápido. Em $N=1.000$, surpresa, Fleury+Naïve foi mais rápido (cerca de 6~ms contra 28~ms). À primeira vista parece contradizer a Tabela~\ref{tab:bridges}, mas é coerente quando se olha o uso real:
\begin{itemize}
    \item o naïve usa \texttt{isReachable(u, v)}, que para a DFS assim que encontra $v$. Em grafos esparsos isso varre poucos vértices;
    \item o Tarjan, mesmo só pra responder se uma aresta é ponte, precisa terminar a DFS inteira porque computa todas as pontes;
    \item nosso cache ajuda quando o Fleury olha vários vizinhos do mesmo vértice no mesmo passo, mas a cada \texttt{removeEdge} a versão muda e o cache é invalidado.
\end{itemize}

Em $N=10.000$ o Fleury+Tarjan gasta cerca de 21 segundos, compatível com a estimativa $E\cdot(V+E)\approx 2\times 10^8$. Em $N=100.000$, qualquer Fleury fica inviável, mas a detecção pura de pontes (Tabela~\ref{tab:bridges}) continua rodando em 20~ms.

\subsection{Por que Tarjan ganha em \texttt{findBridges} mas perde no Fleury?}

A pergunta é diferente em cada caso. Quando se quer enumerar todas as pontes do grafo, o Tarjan resolve em uma DFS e o naïve precisa de $E$ delas. Quando se quer só saber se uma aresta específica é ponte, o naïve termina assim que confirma a alcançabilidade, enquanto o Tarjan tem que terminar a DFS para garantir que está computando \texttt{low} corretamente.

Ou seja, o método certo depende do que você está perguntando. Para listar pontes, Tarjan domina. Para checar uma de cada vez em grafo esparso, o naïve compete bem -- o que foi um resultado contraintuitivo até a gente parar para pensar no porquê.

\section{Como reproduzir}

\begin{lstlisting}[language=bash]
# Compilar
javac *.java

# Modo interativo (menu)
java Main

# Experimentos completos (cuidado: pode demorar em N grande)
java Main --experiments

# Experimentos modo seguro (pula naive em N grande)
java Main --experiments-safe
\end{lstlisting}

Os tempos detalhados são salvos em \texttt{resultados/experimentos.csv}.

\section{Conclusão}

O Tarjan ganha do naïve no que importa para enumerar pontes -- até 108$\times$ em $N=1.000$, e é o único viável em $N=100.000$ (cerca de 20~ms).

No Fleury a comparação é mais sutil. Em grafos esparsos, a terminação antecipada do \texttt{isReachable} no naïve compensa a vantagem assintótica do Tarjan, e o naïve às vezes vence. Em qualquer dos dois, o custo total do Fleury é $O(E\cdot(V+E))$, e isso torna o caso $N=100.000$ inviável -- algoritmos modernos como o de Hierholzer ($O(V+E)$) seriam preferíveis para esse tamanho. Mas o enunciado pediu Fleury, e Fleury foi o que fizemos.

\begin{thebibliography}{9}
\bibitem{tarjan1974}
TARJAN, R. E.
\emph{A note on finding the bridges of a graph}.
Information Processing Letters, v. 2, n. 6, p. 160--161, 1974.
\href{https://doi.org/10.1016/0020-0190(74)90003-9}{https://doi.org/10.1016/0020-0190(74)90003-9}.

\bibitem{cormen}
CORMEN, T. H.; LEISERSON, C. E.; RIVEST, R. L.; STEIN, C.
\emph{Introduction to Algorithms}. 3.~ed. MIT Press, 2009.

\bibitem{euler1736}
EULER, L.
\emph{Solutio problematis ad geometriam situs pertinentis}. Comment. Acad. Sci. U. Petrop., 1741.
\end{thebibliography}

\end{document}
